% =====================================================================
\chapter{Cuestiones}
\label{cap:test}
% =====================================================================

Veinte cuestiones de opción múltiple, de 2 puntos cada una (40 puntos en total), cuatro por módulo. Hay una sola respuesta correcta en cada cuestión y las respuestas erróneas no restan. La clave y su justificación están en el \cref{ap:soluciones}.

\section*{Módulo 24. Formulación y resolución}
\phantomsection\addcontentsline{toc}{section}{Módulo 24. Formulación y resolución}

\begin{cuestiones}
\item \pt{2} Una línea de envasado puede producir como máximo 120~t por semana. Con inventario final nulo, la demanda pendiente desde la semana~$t$ hasta el final del horizonte es de 90~t. En la restricción $x_t\le M\,y_t$, con $y_t\in\{0,1\}$ la variable de preparación, la menor $M$ válida es:
\begin{opciones}
  \item $120$~t;
  \item $210$~t;
  \item $90$~t;
  \item $10^{6}$~t, para no cortar ninguna solución.
\end{opciones}

\item \pt{2} En el óptimo de un programa lineal que maximiza el margen semanal, a la restricción de horas de granuladora le sobran 12~h. Su precio sombra:
\begin{opciones}
  \item es cero;
  \item es el margen óptimo dividido entre 12~h;
  \item coincide con el margen por hora del producto que más granuladora consume;
  \item no puede conocerse sin resolver el dual completo.
\end{opciones}

\item \pt{2} La demanda semanal de un producto perecedero es $\N(200,\,40^{2})$~hL. Cada hectolitro que falta cuesta $c_u=30$~€ y cada hectolitro que sobra, $c_o=10$~€. La cantidad óptima es, aproximadamente:
\begin{opciones}
  \item $200$~hL;
  \item $173$~hL;
  \item $234$~hL;
  \item $227$~hL.
\end{opciones}

\item \pt{2} Sean dos objetivos que se minimizan y pesos $w_1,w_2>0$. Si $x^{\ast}$ minimiza $w_1f_1+w_2f_2$ en el conjunto factible:
\begin{opciones}
  \item $x^{\ast}$ puede estar dominado por otro punto factible;
  \item $x^{\ast}$ es Pareto-óptimo, aunque puede haber puntos del frente que ninguna suma ponderada alcance;
  \item variando los pesos se obtiene siempre el frente de Pareto completo;
  \item $x^{\ast}$ es Pareto-óptimo solo si $f_1$ y $f_2$ son convexas.
\end{opciones}
\end{cuestiones}

\section*{Módulo 25. Optimización bayesiana y modelos sustitutos}
\phantomsection\addcontentsline{toc}{section}{Módulo 25. Optimización bayesiana y modelos sustitutos}

\begin{cuestiones}[resume]
\item \pt{2} En un problema de minimización, el mejor valor observado es $\fmejor=10$ y la posterior en un candidato es $f(x)\sim\N(10,\,2^{2})$. Su mejora esperada vale:
\begin{opciones}
  \item $0$, porque la media no mejora $\fmejor$;
  \item $0{,}399$;
  \item $0{,}798$;
  \item $2$.
\end{opciones}

\item \pt{2} Un proceso gaussiano de media nula con núcleo exponencial cuadrático, $\sigma_f=\ell=1$, se condiciona a una sola observación sin ruido en $x=0$. Varianza posterior de $f(1)$:
\begin{opciones}
  \item $0{,}632$;
  \item $0{,}368$;
  \item $1$, porque la observación no informa sobre otros puntos;
  \item $0{,}135$.
\end{opciones}

\item \pt{2} En una búsqueda aleatoria de hiperparámetros se quiere que, con probabilidad $0{,}95$, al menos una configuración caiga en la región del mejor $5\,\%$. Número mínimo de configuraciones independientes:
\begin{opciones}
  \item $20$;
  \item $19$;
  \item $100$;
  \item $59$.
\end{opciones}

\item \pt{2} Una biblioteca devuelve la desviación típica de una nueva \emph{observación}, que es $\sqrt{\sigma_n^{2}(x)+\sigma_\varepsilon^{2}}$, y por error se usa en la mejora esperada de un problema con ruido. Entonces, en un punto ya ensayado:
\begin{opciones}
  \item la mejora esperada se anula, como debe ser;
  \item la mejora esperada no baja del nivel que fija el ruido y el algoritmo tiende a repetir ensayos casi iguales;
  \item la mejora esperada se vuelve negativa;
  \item no hay efecto apreciable si $\sigma_\varepsilon$ es menor que $\sigma_f$.
\end{opciones}
\end{cuestiones}

\section*{Módulo 26. Control predictivo y aprendizaje por refuerzo}
\phantomsection\addcontentsline{toc}{section}{Módulo 26. Control predictivo y aprendizaje por refuerzo}

\begin{cuestiones}[resume]
\item \pt{2} Un modelo térmico continuo tiene un único autovalor $\lambda=-5$~h$^{-1}$. El modelo discreto obtenido con Euler explícito es estable si y solo si el periodo de muestreo cumple:
\begin{opciones}
  \item $\Delta t<0{,}4$~h;
  \item $\Delta t<0{,}2$~h;
  \item cualquier $\Delta t>0$, porque el sistema continuo es estable;
  \item $\Delta t<0{,}5$~h.
\end{opciones}

\item \pt{2} Sea $x_{k+1}=x_k+u_k$ con coste $x_0^{2}+u_0^{2}+x_1^{2}$ (horizonte $N=1$, $Q=R=P_f=1$) y $x_0=2$. La decisión óptima es:
\begin{opciones}
  \item $u_0=-2$;
  \item $u_0=-0{,}5$;
  \item $u_0=-1$;
  \item $u_0=-1{,}2$.
\end{opciones}

\item \pt{2} En un MPC convexo, la restricción de confort tiene en el óptimo un multiplicador de Karush--Kuhn--Tucker de 3~€ por K·h. Se ablanda con una holgura penalizada con $\rho\sum_j s_j$. La solución blanda coincide con la dura, cuando esta es factible, si:
\begin{opciones}
  \item $\rho>0$;
  \item $\rho\ge 1/3$;
  \item $\rho\le 3$;
  \item $\rho\ge 3$.
\end{opciones}

\item \pt{2} Un episodio de climatización se corta a las 72~h por un límite de tiempo (\texttt{truncated}) sin que el proceso haya terminado. El objetivo del crítico en la última transición $(s,a,r,s')$ debe ser:
\begin{opciones}
  \item $y=r$;
  \item $y=r+\gamma\,Q_{\phi'}(s',\tilde a')$;
  \item $y=\gamma\,Q_{\phi'}(s',\tilde a')$;
  \item $y=r-\gamma\,Q_{\phi'}(s',\tilde a')$.
\end{opciones}
\end{cuestiones}

\section*{Módulo 27. Optimización geométrica y 3D}
\phantomsection\addcontentsline{toc}{section}{Módulo 27. Optimización geométrica y 3D}

\begin{cuestiones}[resume]
\item \pt{2} Con la interpolación SIMP, $p=3$ y $\Emin$ despreciable, un elemento de densidad $x_e=0{,}6$ aporta una rigidez relativa a la del sólido de:
\begin{opciones}
  \item $0{,}6$;
  \item $0{,}36$;
  \item $0{,}216$;
  \item $0{,}4$.
\end{opciones}

\item \pt{2} En el problema de flexibilidad mínima, $c=F\T U$ con $K(x)U=F$ y $N$ densidades. Además de la resolución de $KU=F$, el gradiente $\nabla_xc$ completo necesita:
\begin{opciones}
  \item ninguna resolución adicional, porque el adjunto es $\lambda=U$;
  \item una resolución adjunta con otro término independiente;
  \item $N$ resoluciones, una por densidad;
  \item $2N$ resoluciones por diferencias centradas.
\end{opciones}

\item \pt{2} Una cámara con campo de visión vertical de $60^{\circ}$ produce imágenes de $1080$ píxeles de alto. Un nivel de detalle con error geométrico de $0{,}02$~m puede usarse sin superar $1$~píxel de error a partir de una distancia de:
\begin{opciones}
  \item $10{,}8$~m;
  \item $37{,}4$~m;
  \item $32{,}4$~m;
  \item $18{,}7$~m.
\end{opciones}

\item \pt{2} Una malla triangular cerrada con topología de esfera ($\chi=2$) tiene 1000 vértices. Su número de caras es:
\begin{opciones}
  \item $1996$;
  \item $2000$;
  \item $2994$;
  \item $998$.
\end{opciones}
\end{cuestiones}

\section*{Módulo 28. Búsqueda con LLM, producción y valor}
\phantomsection\addcontentsline{toc}{section}{Módulo 28. Búsqueda con LLM, producción y valor}

\begin{cuestiones}[resume]
\item \pt{2} Cincuenta instrucciones de un clasificador tienen la misma exactitud verdadera y se evalúan con errores $\N(0,\,0{,}04^{2})$. La cota de la maldición del optimizador para el sesgo de la instrucción ganadora es:
\begin{opciones}
  \item $0{,}040$;
  \item $0{,}112$;
  \item $0{,}280$;
  \item $0{,}0057$.
\end{opciones}

\item \pt{2} Una variable binaria de un sensor tenía proporciones $p=(0{,}5;\ 0{,}5)$ en entrenamiento y $q=(0{,}6;\ 0{,}4)$ en la última semana. Su índice de estabilidad poblacional es:
\begin{opciones}
  \item $0{,}0203$;
  \item $0{,}10$;
  \item $0$, porque las medias coinciden;
  \item $0{,}0405$.
\end{opciones}

\item \pt{2} Un halving sucesivo con $\eta=3$ y $K=3$ escalones ($r_{\max}=27\,r_{\min}$, $N$ múltiplo de 27 y coste proporcional al presupuesto) cuesta, respecto a evaluar las $N$ configuraciones con el presupuesto completo, una fracción de:
\begin{opciones}
  \item $1/27$;
  \item $1/3$;
  \item $4/27$;
  \item $1/4$.
\end{opciones}

\item \pt{2} Un optimizador cuesta $I=50\,000$~€ y aporta un flujo neto constante de $20\,000$~€ al final de cada uno de 4~años. Con un tipo de descuento del $10\,\%$, su valor actual neto es, aproximadamente:
\begin{opciones}
  \item $30\,000$~€;
  \item $13\,400$~€;
  \item $63\,400$~€;
  \item $-6\,600$~€.
\end{opciones}
\end{cuestiones}
